\begingroup
\setlength{\eoceAfterSpace}{0mm}
\let\clearpageforsection\relax
\reviewexercisesheader{}

% 35

\eoce{% Replaces gpa_study_hours
\qt{Nama hewan peliharaan\label{seattle_pet_names}}
Kota Seattle, Washington, memiliki portal data terbuka yang
mencakup hewan peliharaan yang terdaftar di kota tersebut.
Untuk setiap hewan peliharaan yang terdaftar, tersedia informasi
tentang nama dan spesiesnya.
Visualisasi berikut memetakan proporsi anjing dengan nama tertentu
terhadap proporsi kucing dengan nama yang sama.
Ditampilkan 20 nama kucing dan anjing yang paling umum.
Garis diagonal pada grafik adalah garis $x = y$;
jika suatu nama terletak pada garis ini, nama tersebut sama populernya
di kalangan anjing maupun kucing.

\noindent\begin{minipage}[c]{0.4\textwidth}
\raggedright\begin{parts}
\item
    Apakah data ini dikumpulkan sebagai bagian dari eksperimen
    atau studi observasional?
\item
    Apa nama anjing yang paling umum? Apa nama kucing yang paling umum?
\item
    Nama apa saja yang lebih umum untuk kucing daripada anjing?
\item
    Apakah hubungan antara kedua variabel positif atau negatif?
    Apa artinya dalam konteks data ini?
\end{parts}\vspace{5mm}
\end{minipage}
\begin{minipage}[c]{0.05\textwidth}
\ 
\end{minipage}
\begin{minipage}[c]{0.53\textwidth}
\begin{center}
\Figures[Ditampilkan diagram pencar dengan setiap titik diberi label nama hewan peliharaan. Sumbu horizontal menyatakan ``Proporsi kucing'' dan membentang dari 0,002 hingga 0,010. Sumbu vertikal menyatakan ``Proporsi anjing'' dan membentang dari 0,002 hingga 0,010. Terdapat pula garis diagonal (y = x), dan hanya dua titik yang berada di bawah garis ini: ``Oliver'' sekitar (0,0045; 0,004) dan ``Lily'' sekitar (0,005; 0,004). Data menunjukkan tren positif yang lemah. Kasus paling ekstrem (proporsi tertinggi untuk anjing atau kucing) adalah ``Lucy'' pada (0,006; 0,0095), ``Charlie'' pada (0,005; 0,009), ``Luna'' pada (0,0065; 0,007), dan ``Bella'' pada (0,005; 0,007).]{0.95}{eoce/seattle_pet_names}{seattle_pet_names}
\end{center}
\end{minipage}
}{}

% 36

\eoce{\qt{Stres, Bagian II\label{stressed_out_experiment}} Dalam suatu studi yang mengevaluasi
hubungan antara stres dan kram otot, separuh subjek ditetapkan secara acak untuk mengalami peningkatan stres dengan ditempatkan di dalam lift yang turun cepat dan berhenti mendadak, sedangkan separuh lainnya dibiarkan tanpa stres atau pada tingkat stres dasar.
\begin{parts}
\item Jenis studi apakah ini?
\item Dapatkah studi ini digunakan untuk menyimpulkan hubungan sebab-akibat antara peningkatan stres
dan kram otot?
\end{parts}
}{}

% 37

\eoce{\qt{Biji chia dan penurunan berat badan\label{chia_weight_lostt}} Chia Pets---patung kecil
terakota yang menumbuhkan ``rambut'' hijau lebat---membuat tanaman chia dikenal luas.
Namun, chia kemudian memperoleh reputasi yang sama sekali baru sebagai suplemen diet.
Dalam sebuah studi tahun 2009, tim peneliti merekrut 38 laki-laki dan membagi mereka
secara acak ke dalam dua kelompok: perlakuan atau kontrol. Mereka juga merekrut 38 perempuan,
lalu menetapkan secara acak separuh peserta ini ke kelompok perlakuan dan separuh lainnya
ke kelompok kontrol. Satu kelompok diberi 25 gram biji chia dua kali sehari, sedangkan
kelompok lain diberi plasebo. Para subjek menjadi sukarelawan untuk mengikuti studi.
Setelah 12 minggu, para ilmuwan tidak menemukan perbedaan signifikan antarkelompok dalam
hal nafsu makan maupun penurunan berat badan.
\footfullcite{Nieman:2009}
\begin{parts}
\item Jenis studi apakah ini?
\item Apa perlakuan eksperimen dan perlakuan kontrol dalam studi ini?
\item Apakah pemblokan digunakan dalam studi ini? Jika ya, apa variabel pemblokannya?
\item Apakah pembutaan digunakan dalam studi ini?
\item Jelaskan apakah kita dapat membuat pernyataan kausal, dan nyatakan apakah
kesimpulannya dapat digeneralisasikan ke populasi yang lebih luas.
\end{parts}
}{}

% 38

\eoce{\qt{Survei dewan kota\label{city_council_survey}}
Sebuah dewan kota meminta agar survei rumah tangga dilakukan
di wilayah pinggiran kotanya.
Wilayah tersebut terbagi menjadi banyak lingkungan yang berbeda-beda:
sebagian berisi rumah besar, sebagian hanya apartemen, dan sebagian lainnya
memiliki campuran struktur hunian yang beragam.
Untuk setiap bagian di bawah ini, tentukan metode pengambilan sampel yang
dijelaskan, lalu uraikan kelebihan dan kekurangan statistik metode tersebut
dalam konteks kota ini.
\begin{parts}
\item
    Ambil sampel acak 200 rumah tangga dari kota tersebut.
\item
    Bagi kota menjadi 20 lingkungan, lalu ambil sampel 10 rumah tangga
    dari setiap lingkungan.
\item
    Bagi kota menjadi 20 lingkungan, ambil sampel acak 3 lingkungan,
    lalu ambil semua rumah tangga dari ketiga lingkungan tersebut sebagai sampel.
\item
    Bagi kota menjadi 20 lingkungan, ambil sampel acak 8 lingkungan,
    lalu ambil sampel acak 50 rumah tangga dari lingkungan-lingkungan tersebut.
\item
    Ambil 200 rumah tangga yang paling dekat dengan kantor dewan kota sebagai sampel.
\end{parts}
}{}

% 39

\eoce{\qt{Penalaran yang cacat\label{flawed_reasoning}} Temukan cacat penalaran
dalam skenario-skenario berikut. Jelaskan apa yang seharusnya dilakukan secara berbeda
oleh orang-orang dalam studi tersebut jika mereka ingin membuat kesimpulan sekuat itu.
\begin{parts}
\item Para siswa di sebuah sekolah dasar diberi kuesioner yang harus dikembalikan
setelah diisi oleh orang tua mereka. Salah satu pertanyaannya berbunyi,
``Apakah jadwal kerja Anda menyulitkan Anda meluangkan waktu bersama anak-anak
sepulang sekolah?'' Dari orang tua yang menjawab, 85\% mengatakan ``tidak''.
Berdasarkan hasil ini, pihak sekolah menyimpulkan bahwa sebagian besar orang tua
tidak mengalami kesulitan meluangkan waktu bersama anak-anak mereka sepulang sekolah.
\item Survei dilakukan terhadap sampel acak sederhana berisi 1.000 perempuan yang
baru saja melahirkan untuk menanyakan apakah mereka merokok selama kehamilan.
Tiga tahun kemudian dilakukan survei tindak lanjut yang menanyakan apakah anak-anak
mereka mengalami gangguan pernapasan. Namun, hanya 567 perempuan yang dapat dihubungi
di alamat yang sama. Peneliti melaporkan bahwa 567 perempuan ini mewakili semua ibu.
\item Seorang dokter ortopedi memberikan kuesioner kepada 30 pasiennya yang tidak
memiliki masalah sendi dan menemukan bahwa 20 di antaranya rutin berlari.
Ia menyimpulkan bahwa berlari menurunkan risiko masalah sendi.
\end{parts}
}{}

% 40

\eoce{\qt{Pendapatan dan pendidikan di county Amerika Serikat\label{income_education_county}}
Diagram pencar di bawah ini menunjukkan hubungan antara pendapatan per kapita
(dalam ribuan dolar) dan persentase penduduk yang memiliki gelar sarjana di
3.143 county di Amerika Serikat pada tahun 2010.

\noindent\begin{minipage}[c]{0.44\textwidth}
\begin{parts}
\item Apa variabel penjelas dan variabel responsnya?
\item Jelaskan hubungan antara kedua variabel. Pastikan untuk membahas pengamatan
yang tidak biasa, jika ada.
\item Dapatkah kita menyimpulkan bahwa memiliki gelar sarjana meningkatkan pendapatan seseorang?
\end{parts}\vspace{8mm}
\end{minipage}
\begin{minipage}[c]{0.55\textwidth}
\begin{center}
\Figures[Ditampilkan diagram pencar dengan ``Persentase bergelar sarjana'' pada sumbu horizontal (0\% hingga 80\%) dan ``Pendapatan per kapita'' pada sumbu vertikal (\$0 hingga \$65.000). Terlihat ribuan titik. Untuk persentase bergelar sarjana antara 0\% dan 20\%, titik-titik biasanya berada pada rentang vertikal \$10.000 hingga \$25.000. Untuk 20\% hingga 40\% pada sumbu horizontal, sebagian besar titik berada antara \$15.000 dan \$35.000 pada sumbu vertikal. Untuk 40\% hingga 60\%, sebagian besar titik berada antara \$25.000 dan \$45.000. Hanya sekitar lima titik yang persentasenya melebihi 60\%, dan semuanya berada di atas \$45.000 pada sumbu vertikal.]{0.78}{eoce/county_income_education}{county_income_education_scatterplot}
\end{center}
\end{minipage}
}{}

% 41

\eoce{\qt[?]{Makan lebih baik, merasa lebih baik\label{eat_better_feel_better}}
Dalam studi kesehatan masyarakat tentang pengaruh konsumsi buah dan sayur terhadap
kesejahteraan psikologis orang dewasa muda, peserta diacak ke tiga kelompok:
(1) pola makan biasa; (2) intervensi ekologis sesaat berupa pengingat pesan teks
untuk meningkatkan konsumsi buah dan sayur beserta kupon untuk membelinya; atau (3) intervensi
buah dan sayur berupa dua porsi tambahan buah dan sayur segar per hari di luar pola
makan biasa. Para peserta adalah mahasiswa sukarelawan di University of Otago,
Selandia Baru, dan mengisi survei melalui ponsel cerdas setiap malam. Pada akhir studi,
hanya kelompok ketiga yang menunjukkan peningkatan kesejahteraan psikologis sepanjang
14 hari dibandingkan dua kelompok lain.\footfullcite{conner2017let}
\begin{parts}
\item
    Jenis studi apakah ini?
\item
    Tentukan variabel penjelas dan variabel responsnya.
\item
    Dapatkah hasil studi digeneralisasikan ke populasi? Jelaskan.
\item
    Dapatkah hasil studi menetapkan hubungan sebab-akibat? Jelaskan.
\item
    Sebuah artikel surat kabar menyatakan,
    ``Hasil studi ini membuktikan bahwa memberi orang dewasa muda buah dan sayuran
    segar untuk dimakan dapat memberikan manfaat psikologis, bahkan dalam waktu singkat.''
    Bagaimana pernyataan itu perlu direvisi agar didukung oleh studi?
\end{parts}
}{}

% 42

\eoce{\qt{Layar, remaja, dan kesejahteraan psikologis\label{screen_time_well_being}}
Studi terhadap tiga kumpulan data berskala besar yang mewakili populasi nasional
Irlandia, Amerika Serikat, dan Britania Raya (n = 17.247) meminta remaja usia 12--15 tahun
mencatat waktu layar dalam buku harian serta menjawab pertanyaan tentang
perasaan atau perilaku; jawaban
itu digunakan untuk menghitung skor kesejahteraan psikologis. Analisis juga menyertakan
jenis kelamin dan usia anak serta, untuk ibu, tingkat pendidikan, etnisitas, tekanan
psikologis, dan status pekerjaan. Studi tersebut menyimpulkan bahwa hanya ada sedikit
bukti tegas bahwa waktu layar menurunkan kesejahteraan remaja.\footfullcite{orben2018screens}
\begin{parts}
\item
    Jenis studi apakah ini?
\item
    Tentukan variabel penjelasnya.
\item
    Tentukan variabel responsnya.
\item
    Dapatkah hasil studi digeneralisasikan ke populasi? Mengapa?
\item
    Dapatkah hasil studi menetapkan hubungan sebab-akibat? Jelaskan.
\end{parts}
}{}

% 43

\eoce{\qt{Stanford Open Policing\label{stanford_open_policing}}
Proyek Stanford Open Policing menghimpun, menganalisis, dan menerbitkan catatan
pemberhentian lalu lintas oleh penegak hukum di seluruh Amerika Serikat untuk membantu
peneliti, jurnalis, dan pembuat kebijakan menyelidiki dan memperbaiki interaksi polisi
dengan masyarakat.\footfullcite{pierson2017large}
Berikut cuplikan tabel ringkasan data proyek tersebut.
\begin{center}
\begin{tabular}{lllrrr}
\hline
               &           & Ras       & Jumlah       & \multicolumn{2}{c}{\% dari yang dihentikan}  \\
County         & Neg. bagian & pengemudi & per tahun     & mobil digeledah & pengemudi ditangkap \\
\hline
Apache County  & Arizona   & Hitam     & 266          & 0.08          & 0.02 \\
Apache County  & Arizona   & Hispanik  & 1008         & 0.05          & 0.02 \\
Apache County  & Arizona   & Putih     & 6322         & 0.02          & 0.01 \\
Cochise County & Arizona   & Hitam     & 1169         & 0.05          & 0.01 \\
Cochise County & Arizona   & Hispanik  & 9453         & 0.04          & 0.01 \\
Cochise County & Arizona   & Putih     & 10826        & 0.02          & 0.01 \\
$\cdots$       & $\cdots$  & $\cdots$  & $\cdots$     & $\cdots$      & $\cdots$ \\
Wood County    & Wisconsin & Hitam     & 16           & 0.24          & 0.10 \\
Wood County    & Wisconsin & Hispanik  & 27           & 0.04          & 0.03 \\
Wood County    & Wisconsin & Putih     & 1157         & 0.03          & 0.03 \\
\hline 
\end{tabular}
\end{center}
\begin{parts}
\item
    Variabel apa yang dikumpulkan untuk setiap pemberhentian lalu lintas
    guna membuat tabel ini?
\item
    Klasifikasikan setiap variabel sebagai numerik (kontinu atau diskret)
    atau kategoris (ordinal atau nonordinal).
\item
    Anda ingin menilai apakah tingkat penggeledahan kendaraan berbeda menurut
    ras pengemudi. Mana variabel respons dan variabel penjelasnya?
\end{parts}
}{}

% 44

\eoce{\qt{Peluncuran antariksa\label{space_launches}}
Tabel berikut merangkum jumlah peluncuran antariksa di Amerika Serikat menurut
jenis lembaga peluncur dan hasilnya (berhasil atau gagal).\footfullcite{data:spacelaunches}
\begin{center}
\begin{tabular}{l | rr | rr}
\hline
        & \multicolumn{2}{| c}{1957 - 1999} & \multicolumn{2}{| c}{2000 - 2018} \\
        & Gagal & Berhasil & Gagal & Berhasil \\
\hline
Swasta  &      13 &     295 &      10 &     562 \\
Negara  &     281 &    3751 &      33 &     711 \\
Rintisan &       - &       - &       5 &      65 \\
\hline
\end{tabular}
\end{center}
\begin{parts}
\item
    Variabel apa yang dikumpulkan untuk setiap peluncuran guna membuat tabel ini?
\item
    Klasifikasikan setiap variabel sebagai numerik (kontinu atau diskret)
    atau kategoris (ordinal atau nonordinal).
\item
    Anda ingin mempelajari bagaimana tingkat keberhasilan berbeda menurut
    lembaga peluncur dan periode waktu. Mana variabel respons dan penjelasnya?
\end{parts}
}{}
\endgroup
